\documentclass[10pt,a4paper]{amsart}
\usepackage[margin=19mm]{geometry}
\usepackage{amsmath,amssymb,mathtools}
\usepackage[scr=rsfs,cal=euler]{mathalfa}
\usepackage{xfrac}
\usepackage[unicode,hidelinks,bookmarks=false]{hyperref}
\newcommand{\ie}{i.e.\ }
\newcommand{\cf}{cf.\ }
\newcommand{\resp}{respectively\ }
\newcommand{\Subsection}{\subsection}
\providecommand{\nobreakdash}{\nobreak\hbox{-}}
\setlength{\emergencystretch}{2em}
\multlinegap0pt
\usepackage{bm}
\usepackage{bbm}
\usepackage{dsfont}
\usepackage{xspace}
\usepackage{cancel}
\usepackage{mathtools}
\usepackage{amsthm}
\makeatletter
\@ifundefined{lemma}{\newtheorem{lemma}{Lemma}}{}
\@ifundefined{remark}{\newtheorem{remark}{Remark}}{}
\makeatother
\title[Global solutions for supersonic flow of a Chaplygin gas past]{Global solutions for supersonic flow of a Chaplygin gas past a conical wing with a shock wave detached from the leading edges}
\author{Bingsong Long}
\date{Source: arXiv:2503.03406v4 (CC BY 4.0)}
\begin{document}
\maketitle

\noindent\emph{Source and licence.} Global solutions for supersonic flow of a Chaplygin gas past a conical wing with a shock wave detached from the leading edges. Version arXiv:2503.03406v4 (CC BY 4.0), distributed under the Creative Commons Attribution 4.0 International licence, as recorded in the arXiv licence metadata for this exact version and shown on the abstract page https://arxiv.org/abs/2503.03406v4. Full rights notice: licenses/arxiv-2503.03406-CC-BY-4.0.txt.\par\smallskip

\noindent\textit{Editorial scope.} Counted unit: the Lemma whose source label is \texttt{lemma: inf-estimate} in the article (source0.tex lines 462--472, source label \texttt{lemma: inf-estimate}), delivered together with the article's own complete original proof of it (source0.tex lines 474--554, one \texttt{proof} environment of 81 lines). No mathematics is written, repaired or re-derived: statement and proof are the article's own text line for line. Required context retained uncounted: eq: conical tran. (lines 259--261); eq: for Fmu (lines 363--365); eq: spherical tran. (lines 402--404); eq: psi zhankai (lines 410--413); eq:nianxing BC4 (lines 415--422); lemma: comparison principle (lines 427--435); eq:3.42 (lines 574--576); remark:lip on supsolution (lines 607--609). Every definition, display and label of those lines is retained unchanged. Omitted from this package and not redistributed: 1--258, 262--362, 366--401, 405--409, 414--414, 423--426, 436--461, 473--473, 555--573, 577--606, 610--874. Nothing inside a retained excerpt is omitted or altered: every retained body is delivered exactly as the article's lines are written, so no omission receipt of any kind accompanies this package. Citations: the retained excerpts carry no citation command at all, so no bibliography entry of the article is transcribed and no reference is redistributed. Reference adaptation: none was needed and none was made; every reference inside the delivered unit resolves inside it, so no unresolved reference or citation is produced. Printed slips kept literal: no printed word, symbol, hypothesis or number of the counted statement or of its proof was altered. Layout and class: the article's own class is \texttt{amsart} with packages including ragged2e, amssymb, amsfonts, amsmath, amsthm, mathrsfs, diagbox, color, soul, bbm, graphicx, empheq, framed, bm; the wrapper is the lane's pinned \texttt{amsart} editorial wrapper at 10pt on A4, which restates the theorem-like environments the retained text needs and adds the article's own macro definitions that the retained text needs (none), with extra wrapper packages (bm, bbm, dsfont, xspace, cancel, mathtools). The delivered numbering is the unit's own. Assets: no retained span contains a figure environment, a graphics-inclusion command, a PDF-page inclusion, a table or a TikZ picture; no image file is copied into this package, the package declares no asset dependency and no image file is redistributed with it. Licence: version arXiv:2503.03406v4 of this article is distributed under the Creative Commons Attribution 4.0 International licence, as recorded in the arXiv licence metadata for this exact version and shown on the abstract page https://arxiv.org/abs/2503.03406v4. A full rights notice is delivered at licenses/arxiv-2503.03406-CC-BY-4.0.txt. Characters: every retained excerpt is pure ASCII, so the delivered engine, XeLaTeX with its default Latin Modern text font, reports no missing character.
\begin{equation}\label{eq: conical tran.} 
	\rho(\bm{x})=\rho(\bm{\xi}), \quad \Phi(\bm{x})=x_3\phi(\bm{\xi}),
\end{equation}

\begin{equation}\label{eq: for Fmu}
	\mathcal{G}(\mu,\phi):=c^{2}(\Delta\phi+D^{2}\phi[\bm{\xi},\bm{\xi}])-\mu D^{2}\phi[D\phi-\chi\bm{\xi},D\phi-\chi\bm{\xi}]=0\quad\text{in $\Omega$}
\end{equation}

\begin{align}\label{eq: spherical tran.}
    \rho(\bm{x})=\rho(\bm{\zeta}), \quad \Phi(\bm{x})=r\psi(\bm{\zeta}),
\end{align}

\begin{multline}\label{eq: psi zhankai}
    c^{2}\big(\dfrac{\partial_{\theta}(\sin\theta\partial_{\theta}\psi)}{\sin\theta}+\dfrac{\partial_{\varphi\varphi} \psi}{\sin^2\theta}+2\psi\big)-\mu\Big(\partial^{2}_\theta\psi\partial_{\theta\theta}\psi+\dfrac{\partial^2_\varphi\psi}{\sin^2\theta}\dfrac{\partial_{\varphi\varphi} \psi}{\sin^2\theta}\\+2\partial_\theta \psi\dfrac{ \partial_\varphi\psi}{\sin\theta}\dfrac{\partial_{\theta\varphi}\psi}{\sin\theta}-\cot\theta\dfrac{\partial^2_\varphi\psi}{\sin^2\theta}\partial_\theta\psi+\big(\partial^{2}_\theta\psi+\dfrac{\partial^2_\varphi\psi}{\sin^2\theta}\big)\psi\Big)
    =0\quad\text{in $\Omega$}
\end{multline}

  \begin{equation}\label{eq:nianxing BC4}
		\begin{cases}
	        \psi=1+\varepsilon \quad&\text{on $\Gamma_{cone}^{\infty}$},\\ 
			D_{\bm{\zeta}}\psi\cdot\bm{\nu}_{py}=0\quad&\text{on $\Gamma_{py}$},\\
			D_{\bm{\zeta}}\psi\cdot\bm{\nu}_{sy1}=0\quad&\text{on $\Gamma_{sy1}$},\\
                D_{\bm{\zeta}}\psi\cdot\bm{\nu}_{sy2}=0\quad&\text{on $\Gamma_{sy2}$},	
		\end{cases}
	\end{equation}

\begin{lemma}\label{lemma: comparison principle}
Let $\Omega_{D}$ be an open bounded domain in the unit sphere that does not contain $\{\theta=0\}$, with piecewise smooth boundary $\partial\Omega_D$ composed of a finite number of curves $\Gamma_i$, where $i=1,\cdots,n$. Also assume that if $\Gamma_k$ and $\Gamma_l$ meet, the formed angle lies in the interval $(0,\pi]$. Suppose that ${\psi}_{\pm}\in C^{0}(\overline{\Omega_{D}})\cap C^{1}(\overline{\Omega_{D}}\setminus\overline{\Gamma_1})\cap C^2(\Omega_{D})$ satisfy ${\psi}_{\pm}>1$ in $\overline{\Omega_{D}}\setminus\overline{\Gamma_1}$. Assume further that for any $\mu\in [0,1]$, the operator $\mathcal{N}_\mu$ is locally uniformly elliptic with respect to either ${\psi}_{+}$ or $\psi_-$, and
	\begin{align*}
	\mathcal{N}_\mu\psi_-\geq0, \quad\mathcal{N}_\mu\psi_+\leq0\quad&\text{in}~\Omega_{D},\\
   \psi_-\leq\psi_+\quad&\text{on}~\Gamma_1,\\
     D_{\bm{\zeta}}\psi_-\cdot\bm{\nu}_{d_j}<D_{\bm{\zeta}}\psi_+\cdot\bm{\nu}_{d_j}\quad&\text{on}~\Gamma_j,
	\end{align*}
where $\bm{\nu}_{d_j}$~is the exterior normal to $\Gamma_j$ with $j=2,\cdots,n$. Then, it follows that ${\psi}_{-}\leq{\psi}_{+}$ in $\overline{\Omega_{D}}\setminus\overline{\Gamma_1}$.
\end{lemma}

	\begin{equation}\label{eq:3.42} 
		\|D{\phi}\|_{L^{\infty}(\Omega)}\leq C.
	\end{equation}

\begin{remark}\label{remark:lip on supsolution}
 Note that the establishment of the Lipschitz estimate \eqref{eq:3.42} depends only on the super-solution $\psi^+$ and is independent of the sub-solution $\psi^-$ obtained in Lemma $\ref{lemma: inf-estimate}$.
\end{remark}

\begin{lemma}\label{lemma: inf-estimate}
	Let $\psi \in C^{0}(\overline{\Omega})\cap C^{1}(\overline{\Omega}\setminus\overline{\Gamma_{cone}^{\infty}})\cap C^2\big(\overline{\Omega}\setminus(\overline{\Gamma_{cone}^{\infty}}\cup\{P_3\}\cup\{P_4\})\big)$ satisfy \eqref{eq: psi zhankai}--\eqref{eq:nianxing BC4} and $\psi\geq 1+\varepsilon$ in $\overline{\Omega}\setminus\overline{\Gamma_{cone}^{\infty}}$. Then there exists a positive constant $C$, which is independent of $\mu$ and $\varepsilon$, such that
	\begin{equation}\label{eq1:3.13}
		1+\varepsilon<{\psi}\leq C\quad\text{in $\overline{\Omega}\setminus\overline{\Gamma_{cone}^{\infty}}$}.
	\end{equation}
    Moreover, for any subdomain $\Omega_{sub}$ that is strictly away from the degenerate boundary $\Gamma_{cone}^{\infty}$, we have
        \begin{equation}\label{eq:strictly bounded}
		1+\varepsilon+\delta_0<{\psi}\leq C,
	\end{equation}
    where the constant $\delta_0>0$ depends only on the domain $\Omega$. 
\end{lemma}

\begin{proof}
	 First, we find an exact solution to equation \eqref{eq: psi zhankai}. From \eqref{eq: conical tran.} and \eqref{eq: for Fmu}, the corresponding $\Phi$ satisfies
	\begin{equation}\label{eq:3.25}
		(|\nabla_{\bm{x}}\Phi|^2-1)\Delta_{\bm{x}}\Phi-\mu D^2_{\bm{x}}\Phi[\nabla_{\bm{x}}\Phi,\nabla_{\bm{x}}\Phi]=0.
	\end{equation}
     Clearly, for any constant vector $\bm{\eta}\in \mathbb{R}^{3}$, the linear function $\Phi^{\bm{\eta}}=\bm{\eta}\cdot\bm{x}$ is a solution to \eqref{eq:3.25}. We deduce from \eqref{eq: spherical tran.} that
	\begin{equation}\label{eq:3.26}
		{\psi}^{\bm{\eta}}(\bm{\zeta})=\bm{\eta}\cdot\frac{\bm{x}}{|\bm{x}|}
	\end{equation}
	is an exact solution to equation \eqref{eq: psi zhankai}. In addition, since
	\begin{align}\label{eq: d phi first}
            \begin{split}
	    \partial_\theta {\psi}^{\bm{\eta}}(\bm{\zeta})&=-\eta_1\sin\theta+\eta_2\cos\theta\cos\varphi+\eta_3\cos\theta\sin\varphi,\\
            \dfrac{\partial_\varphi {\psi}^{\bm{\eta}}(\bm{\zeta})}{\sin\theta}&=-\eta_2 \sin\varphi+\eta_3\cos\varphi,
            \end{split}
	\end{align}	
    the function ${\psi}^{\bm{\eta}}$ is Lipschitz bounded in the domain $\Omega$.
    
    Next, we construct the super-solutions to problem \eqref{eq: psi zhankai}--\eqref{eq:nianxing BC4}. Let $P^{\bm{\eta}}:=\frac{\bm{\eta}}{|\bm{\eta}|}$ denote the intersection of the vector $\bm{\eta}$ with the unit sphere, and let
    \begin{equation*}
         \Lambda:=\{(\theta,\varphi):\theta\in(0,\frac{\pi}{2}),\varphi\in(\frac{3\pi}{2},2\pi)\}.
    \end{equation*}
    Consider the following set with a fixed $\varepsilon>0$
	\begin{equation}\label{eq:4.8}
		\Sigma_{+}:=\{\bm{\eta}\in\mathbb{R}^{3}: P^{\bm{\eta}}\in \Lambda, \text{ and }{\psi}^{\bm{\eta}}>1+\varepsilon\text{ on }\Gamma_{cone}^{\infty}\}.
	\end{equation}
    Note that for any fixed $\bm{\eta}\in \mathbb{R}^{3}$, the function ${\psi}^{\bm{\eta}}$ is monotonically decreasing in the angle between $\frac{\bm{x}}{|\bm{x}|}$ and $\bm{\eta}$. Hence, when $P^{\bm{\eta}}\in \Lambda$, the values of $D_{\bm{\zeta}}\psi\cdot\bm{\nu}_{py}$ on $\Gamma^{\infty}_{cone}$, $D_{\bm{\zeta}}\psi\cdot\bm{\nu}_{sy1}$ on $\Gamma_{sy1}$, and $D_{\bm{\zeta}}\psi\cdot\bm{\nu}_{sy2}$ on $\Gamma_{sy2}$ are strictly positive. This indicates that
   when $\bm{\eta}\in \Sigma_{+}$, the function ${\psi}^{\bm{\eta}}$ satisfies
	\begin{equation*}
		\begin{cases}
			\mathcal{N}_\mu {\psi}^{\bm{\eta}}= 0\quad&\text{in $\Omega$},\\
			{\psi}^{\bm{\eta}}>1+\varepsilon\quad&\text{on $\Gamma_{cone}^{\infty}$},\\
            D_{\bm{\zeta}}\psi\cdot\bm{\nu}_{py}>0\quad&\text{on $\Gamma_{py}$},\\
			D_{\bm{\zeta}}\psi\cdot\bm{\nu}_{sy1}>0\quad&\text{on $\Gamma_{sy1}$},\\
                D_{\bm{\zeta}}\psi\cdot\bm{\nu}_{sy2}>0\quad&\text{on $\Gamma_{sy2}$}.
		\end{cases}
	\end{equation*}
   Additionally, it follows from \eqref{eq:3.26} that ${\psi}^{\bm{\eta}}>1+\varepsilon$ holds in $\overline{\Omega}\setminus\overline{\Gamma_{cone}^{\infty}}$. Combining this fact and \eqref{eq: d phi first}, the operator $\mathcal{N}_\mu $ is locally uniformly elliptic with respect to ${\psi}^{\bm{\eta}}$. Applying Lemma \ref{lemma: comparison principle}, we conclude that $\psi\leq {\psi}^{\bm{\eta}}$ in $\overline{\Omega}\setminus\overline{\Gamma_{cone}^{\infty}}$. Define $\psi^+$ as the infimum of all such super-solutions. Then, $\psi\leq \psi^+$ holds in $\overline{\Omega}\setminus\overline{\Gamma_{cone}^{\infty}}$. Moreover, thanks to the convexity of $\Gamma_{cone}^{\infty}$, we obtain
	\begin{equation*}
		\psi^+=1+\varepsilon\quad\text{on $\Gamma_{cone}^{\infty}$}.
	\end{equation*}
  
   Finally, we construct the sub-solutions to problem \eqref{eq: psi zhankai}--\eqref{eq:nianxing BC4}. Consider the following set with a fixed $\varepsilon>0$
	\begin{equation}\label{eq: Sigma fu}
		\Sigma_{-}:=\{\bm{\eta}\in\mathbb{R}^{3}: P^{\bm{\eta}}\in \Omega,\text{ and }{\psi}^{\bm{\eta}}<1+\varepsilon~\text{on}~\Gamma_{cone}^{\infty}\},
	\end{equation}
  Similarly, when $\bm{\eta}\in {\Sigma}_{-}$, the function ${\psi}^{\bm{\eta}}$ satisfies
	\begin{equation*}
		\begin{cases}
			\mathcal{N}_\mu {\psi}^{\bm{\eta}}= 0\quad&\text{in $\Omega$},\\
			{\psi}^{\bm{\eta}}<1+\varepsilon\quad&\text{on $\Gamma_{cone}^{\infty}$},\\
            D_{\bm{\zeta}}\psi\cdot\bm{\nu}_{py}<0\quad&\text{on $\Gamma_{py}$},\\
			D_{\bm{\zeta}}\psi\cdot\bm{\nu}_{sy1}<0\quad&\text{on $\Gamma_{sy1}$},\\
                D_{\bm{\zeta}}\psi\cdot\bm{\nu}_{sy2}<0\quad&\text{on $\Gamma_{sy2}$}.
		\end{cases}
	\end{equation*}
  From Remark \ref{remark:lip on supsolution} below, and using the super-solution $\psi^+$ above and the assumption of $\psi\geq 1+\varepsilon$ in $\overline{\Omega}\setminus\overline{\Gamma_{cone}^{\infty}}$, we deduce the locally uniform ellipticity of $\mathcal{N}_\mu$ with respect to $\psi$. Thus, using Lemma \ref{lemma: comparison principle}, we have
	\begin{equation}\label{eq1:11}
		{\psi}\geq {\psi}^{\bm{\eta}}\quad\text{in $\overline{\Omega}\setminus\overline{\Gamma_{cone}^{\infty}}$},
	\end{equation}
  for all $\bm{\eta}\in {\Sigma}_{-}$.  
   
    Define $\psi^-$ as the supremum of all such sub-solutions. We assert here that for any subdomain $\Omega_{sub}$ strictly away from the degenerate boundary $\Gamma_{cone}^{\infty}$, there exists a positive constant $\delta_0$, depending only on $\Omega_{sub}$, such that
	\begin{equation}\label{eq1:3.21}
		{\psi}^{-}\geq 1+\varepsilon+\delta_0\quad\text{in $\Omega_{sub}$}.
	\end{equation}
   Given any point~$\bm{x}_{0}=(x_{10},x_{20},x_{30})$ satisfying $\frac{(x_{10},x_{20},x_{30})}{r}\in \Omega_{sub}$, there exists a sufficiently small constant $0<\delta_0\ll 1$ such that the constant vector $\bm{\eta}_{0}=(1+\varepsilon+\delta_0)\frac{(x_{10},x_{20},x_{30})}{r}\in\Sigma_{-}$. The assertion then follows directly from the definition of $\psi^-$.
 
   Thanks to the continuity of ${\psi}^{-}$, the inequality \eqref{eq1:3.21} means ${\psi}^{-}>1+\varepsilon$ in $\overline{\Omega}\setminus\overline{\Gamma_{cone}^{\infty}}$ and ${\psi}^{-}\geq 1+\varepsilon$ on $\Gamma_{cone}^{\infty}$. Besides, we observe that ${\psi}^{-}\leq 1+\varepsilon$ on $\Gamma_{cone}^{\infty}$. This leads to
	\begin{equation}\label{eq:3.31}
		{\psi}^{-}= 1+\varepsilon\quad\text{on $\Gamma_{cone}^{\infty}$}.
	\end{equation}
 
  In conclusion, we get
	\begin{align}
		1+\varepsilon<{\psi}^-\leq \psi\leq \psi^+\quad&\text{in}~  \overline{\Omega}\setminus\overline{\Gamma_{cone}^{\infty}}\label{eq1:4.19}
		\shortintertext{and}
		\psi=\psi^\pm=1+\varepsilon\quad&\text{on}~ \Gamma_{cone}^{\infty}.\label{eq1:4.20}
	\end{align}
	The proof is complete. 
\end{proof}

\end{document}
